% ============================================================================
% 4 METHODOLOGY
% ============================================================================
\section{Methodology}
\label{sec:methodology}

\paragraph{Experimental control and provenance.}
All reported LLM runs use the configured \code{hosted:gpt-5.4} deployment at
temperature $0$. Each prediction record stores the example and document IDs,
model and method, raw and parsed output, prompt and manifest hashes, token
usage, elapsed time and validation state. Evaluation is a separate deterministic
step, so alternative matching policies can be applied without another API call.
We report only conditions backed by retained prediction and score artifacts.

\paragraph{Short-input baselines.}
Ticket classification predicts one \labelname{type}, one \labelname{queue} and
a set of \labelname{tags}. A zero-shot call selects from allowed values and
quotes evidence; an agentic two-step method first extracts evidence and then
maps it to labels. Both completed the same 1{,}000-ticket slice. Few-shot
classification is implemented but only smoke-tested on one row and is excluded
from comparison. Embedding-based retrieval of reviewed tickets is also
implemented as an optional dynamic few-shot method, but no leakage-safe
reference/development/test split was available for the reported synthetic-ticket
experiment; consequently, it is not included in the reported comparison. For
interactive use, the annotation application additionally implements a pure
nearest-neighbor router: it embeds a ticket locally and applies a
similarity-weighted vote over reviewed reference tickets, without making a
routing LLM call. This route is likewise excluded from the reported accuracy
results because it has not yet been tested on an independent target-domain
split. For Few-NERD, the two directly comparable 1{,}000-sentence
conditions are few-shot structured JSON and few-shot free-form output, each
using three training examples and explicit inclusive token indices. The
zero-shot extraction artifacts contain 100 examples and are reported only as a
supplementary scale check, not as a full-slice comparison.

\paragraph{SciREX prompt selection and frozen evaluation.}
The application derives a deterministic prompt from the active label schema,
definitions and domain guidance; this suggestion costs no LLM call and remains
editable. On a fixed 20-paper development sample, we compared the previous
static few-shot prompt with this suggestion using the first ten app-detected
sentences per window. The selected prompt presents Method, Task, Metric and
Dataset to users; Dataset is mapped to SciREX's release label Material during
evaluation. After prompt selection, a fresh 20-paper official-test sample,
excluding test papers already inspected, was processed in full. A separate
100-paper training-split gate, balanced across four length buckets, tested full
coverage, throughput and failure recovery. No test result was used to revise
the reported prompt.

\paragraph{Document-scale execution.}
Windows are split into sentences only for API calls; all sentences are processed
when \code{max\_sentences=0}. Four workers handle independent documents. Before
each call and after each successful sentence, state is written atomically.
Repeating an identical command resumes at the next unfinished sentence, while
hash checks reject mixed manifests, models or prompts. The runner enforces call
and request-rate ceilings, retries transient failures, and records missing usage
and invalid responses. Provider prices were unavailable, so monetary cost is
not inferred from token counts.

\paragraph{Metrics.}
For extraction, strict one-to-one matching requires the same case-insensitive
label and exact half-open character span. Boundary-tolerant matching permits a
start within one token and an end within two; overlap matching accepts any
same-label character overlap. Strict is primary; the other two quantify useful
localization for human review. We report micro precision, recall and $F_1$,
macro-document $F_1$, per-label and length-bucket results. Uncertainty intervals
use 2{,}000 percentile-bootstrap samples clustered by source \code{doc\_id}.
Ticket single-label fields use accuracy and macro-$F_1$; multi-label tags use
micro-$F_1$. Invalid JSON, invalid labels, evidence validity, failures, token
usage and coverage are reported as operational measures.
